\section{An exact decomposition of the prime sum}
\label{sec:decomposition}

All factorizations in this section are factorizations of ideals prime to
$3$, written using their primary generators.  In particular, prime
factors of the same norm are still regarded as distinct when their
ideals are distinct.  Write $L=\log X$, and let $V$ be a fixed smooth
function supported in a compact subinterval of $(0,\infty)$.
We consider sums with the envelope $V(N(n)/X)$.  Constants below may
depend on its support and on fixed smoothness bounds.

Fix sufficiently small positive constants $\kappa,\rho$, with
$\rho$ sufficiently small in terms of $\kappa$, and then fix
\[
 0<\xi<\kappa/10,\qquad w=X^\xi,\qquad z=X^{.40}.
\]
The choice of $\kappa$ relative to the power-width constant in
Proposition~\ref{prop:dispersion} will be made in the next section.
We call a height block \emph{ordinary} when $T\ll X^{.01}$, including
bounded heights, and \emph{upper} when $T\gg X^{.01}$.  The upper
blocks used below have $T\ll X^{1/6+\rho}$.

The construction reflects the two coefficient regimes in dispersion.
For bilinear norm lengths $AB\asymp X$, with $B\ll\sqrt X$,
the range $B>X^{.40}$ requires independently weighted smooth prime
convolutions; the range $B\le X^{.40}$ permits sharp restrictions on
divisor-bounded coefficients.  We first expose a bounded tuple of
primes with norms between $w$ and $2z$.  If its product scale is large,
we expand the remaining factors into bounded prime tuples and group
their fixed scales.  If that scale is small, we use Mobius inversion
and stop a product of small primes in norm bins.  This second
construction must separate the two coefficients exactly.  On ordinary
blocks it must also retain roughness and \textup{(SW)}, apart from a
sparse family whose contribution can be bounded directly.

\begin{proposition}[Prime decomposition]\label{prop:decomposition}
For either kind of height block, the prime sum with envelope $V$ has
an exact decomposition into a number of Type~I and bilinear terms
bounded by a fixed power of $L$.  This decomposition is valid with
either kernel
\[
 G(n)N(n)^{it}
 \quad\hbox{or}\quad
 c_*\mu^2(n)N(n)^{-1/6+it},
\]
and uses the same coefficients for the two kernels.  The coefficients
are independent of $t$, divisor-bounded, and have all fixed absolute
moments of the usual size $Y L^{O(1)}$ on a norm dyad of length $Y$.
The resulting terms have the following properties.
\begin{enumerate}
\item Type~I terms have an unrestricted inner variable $u$ and an outer
variable $r$ with
\[
 N(r)\asymp R,\qquad
 R\le X^{.40}\quad\hbox{on ordinary blocks},\qquad
 R\le X^{1/3-\kappa/2}\quad\hbox{on upper blocks}.
\]
Their inner weight is a smooth envelope at length $X/R$, with fixed
derivative bounds, depending permissibly on $r$.
\item Bilinear terms have independent coefficients $\alpha_a,\beta_b$,
may be restricted individually to squarefree $a,b$, and satisfy
\[
 N(a)\asymp A,\quad N(b)\asymp B,\quad AB\asymp X,\quad B\ll\sqrt X.
\]
On ordinary blocks, $B\ge cX^{1/3}$ for some fixed $c>0$; on upper
blocks, $B\ge X^{1/3-3\kappa}$.  Whenever $B>X^{.40}$, the $b$
coefficient is a squarefree-restricted convolution of a bounded number
of prime sequences on independent smooth norm dyads, each prime
having norm at least a fixed positive power of $X$.  There is no
additional cutoff coupling those prime dyads.
\item On ordinary blocks, apart from the sparse terms in part~(5),
the $b$ coefficient is supported on numbers all of whose prime factors
have norm greater than $L^{C_r}$ and satisfies the property
\textup{(SW)} required in Proposition~\ref{prop:dispersion}.
Here $C_r$ is any fixed constant chosen sufficiently large later.
\item For any fixed $C_g$, the ordinary bilinear terms with
\[
 1\ll H:=B^2/A<L^{C_g}
\]
consist of precisely three primes: one prime on the $b$ side and two
specified prime dyads on the $a$ side.  All three scales are
$X^{1/3}L^{O(1)}$, their weights and multiplicities are bounded, and
the number of these terms is $O((1+\log L)^{O(1)})$.  In each such
term,
\begin{equation}\label{eq:decomposition-triple-moments}
 \sum_a|\alpha_a|^2\ll A/L^2,
 \qquad
 \sum_b|\beta_b|^2+\sum_b|\beta_b|\ll B/L.
\end{equation}
All other nonsparse ordinary bilinear terms have $H\ge L^{C_g}$
for sufficiently large $X$.
\item The remaining ordinary terms have $X^{.35}\ll B<X^{.40}$
and a sparse $b$ support.  Their total contribution, for either kernel
and uniformly in $t$, is $O(X^{5/6-\eta})$ for some $\eta>0$ depending
on the fixed parameters.
\end{enumerate}
The exponents in the number of terms and in the divisor bounds can be
chosen independently of $C_r$ and $C_g$.  All asserted identities
hold before estimating either kernel.
\end{proposition}

\subsection{Prime detection and bounded prime tuples}

Take a fixed smooth function $\psi$ with values in $[0,1]$, equal to
$1$ on $(0,1]$ and to $0$ on $[2,\infty)$, and put
\[
 \psi_y(\pi)=\psi(N(\pi)/y),\qquad
 \mathcal F_y(n)=\prod_{\pi\mid n}(1-\psi_y(\pi)).
\]
Both kernels in Proposition~\ref{prop:decomposition} vanish on
nonsquarefree arguments.  We may therefore use identities only on
squarefree $n$ throughout this section.

\begin{lemma}[Exact prime detector]\label{lem:prime-detector}
For squarefree $n$ on the envelope support and sufficiently large $X$,
\begin{equation}\label{eq:prime-detector}
 \boldsymbol 1_{\{n\ {\mathrm{prime}}\}}
 =\mathcal F_z(n)-
  \boldsymbol 1_{\{\Omega(n)=2\}}\mathcal F_z(n).
\end{equation}
Moreover, $\mathcal F_z(n)$ is an exact sum of bounded-length
distinguished prime tuples
\begin{equation}\label{eq:distinguished-tuples}
 \mathcal F_z(n)=
 \sum_{i\ge0}\frac1{i!}
 \sum_{n=\pi_1\cdots\pi_i h}
 \mathcal F_w(h)\prod_{j=1}^i
       (\psi_w(\pi_j)-\psi_z(\pi_j)).
\end{equation}
Only boundedly many $i$ contribute, and the distinguished primes
have norms in $[w,2z]$.
\end{lemma}

\begin{proof}
Nonvanishing of $\mathcal F_z(n)$ requires every prime factor to have
norm greater than $z$.  Three factors would give $N(n)>X^{1.20}$,
outside the fixed envelope.  On a prime in that envelope the factor
$\mathcal F_z$ is exactly $1$.  This proves
\eqref{eq:prime-detector}.

Expand, independently at every prime factor,
\[
 1-\psi_z=(1-\psi_w)+(\psi_w-\psi_z).
\]
The chosen factors in the second summand form an unordered subset;
ordering them and dividing by $i!$ gives
\eqref{eq:distinguished-tuples}.  The difference of cutoffs vanishes
below $w$ and above $2z$.  Since the total norm is $O(X)$ and every
chosen prime has norm at least $X^\xi$, their number is bounded in
terms of $\xi$.  The same bound applies to the number of prime factors
of $h$ when $\mathcal F_w(h)\ne0$.
\end{proof}

The two-prime correction in \eqref{eq:prime-detector} is an ordered
pair sum with scalar $1/2$.  Smooth dyadic partitions place its two
prime norms between constant multiples of $X^{.40}$ and $X^{.60}$.
Taking the shorter scale for $B$ gives all the required bilinear
properties, with $H$ a fixed positive power of $X$.

In \eqref{eq:distinguished-tuples}, put $r=\pi_1\cdots\pi_i$ and
insert smooth norm dyads on the distinguished primes.  Let $R$ be
the product of their scales, with $R=1$ when $i=0$.  Thus $N(r)\asymp R$,
with a fixed support ratio.  Set
\[
 s_0=.345\quad\hbox{on ordinary blocks},\qquad
 s_0=1/3-2\kappa\quad\hbox{on upper blocks}.
\]

Suppose first that $R\ge X^{s_0}$.  Expand $h$ by its bounded number
of prime factors, divide the ordered factorization by its factorial,
and insert independent smooth prime dyads.  All prime norms are at
least $w$ and at most a constant times $X^{1-s_0}$.  The distinguished
ones also have norm at most $2z$.  A configuration of scales can be
discarded when it is disjoint from the envelope, but no product or
partial-product cutoff is inserted into its coefficients.

\begin{lemma}[Grouping prime factors]\label{lem:prime-grouping}
The prime factors just obtained can be grouped into scales $A,B$,
with $AB\asymp X$ and $B\ll\sqrt X$, as follows.  On ordinary blocks
either $B\ge X^{1/3+\delta_1}$ for some fixed $\delta_1>0$, or there
are exactly three prime factors with exponents close to $1/3$, and
$B$ is the largest of their three scales.  On upper blocks one can
always arrange $B\ge X^{1/3-3\kappa}$.
\end{lemma}

\begin{proof}
Normalize the logarithms of the prime scales by the logarithm of their
product.  In any limiting configuration they are positive, sum to $1$,
and are bounded below in terms of $\xi$.  On ordinary blocks no part
reaches $2/3$, since every part is at most $1-.345=.655$.

We claim that a positive partition of $1$, with every part less than
$2/3$, has a subset sum in $(1/3,2/3)$ unless it is
$(1/3,1/3,1/3)$.  If there is no such subset, every individual part
is at most $1/3$.  If a part $a<1/3$ exists, start with it and add
parts until the sum first exceeds $1/3$.  The resulting sum is at
most $2/3$.  Equality would require both the preceding sum and the
last added part to be $1/3$; but then $a$ together with that last part
is a subset sum in $(1/3,2/3)$.  Thus equality is also impossible.
The only remaining partition has all parts equal to $1/3$.

There are finitely many possible factor counts.  On the compact sets
of their allowed exponent vectors, outside a fixed small neighborhood
of the exceptional three-part vector, an interior subset has a
uniform distance from both endpoints.  Group that subset and its
complement and call the shorter group $B$.  The fixed support ratios
can be absorbed by reducing the uniform gap, giving the claimed
$\delta_1$.  In the exceptional neighborhood, take the largest prime
scale for $B$; make the neighborhood small enough that $B<X^{.40}$
and $B\ll\sqrt X$.

For upper blocks, a part exceeding $2/3$ can only come from $h$,
and its complement has limiting exponent at least $s_0$.
If no part exceeds $2/3$, a subset sum in $[1/3,2/3]$ exists by the
same successive-sum argument.  Fixed support ratios and limiting
errors are absorbed by the spare $\kappa$ in the asserted lower bound.
\end{proof}

For three factors, their largest scale satisfies $B\ge cX^{1/3}$.
If $H=B^2/A\asymp B^3/X<L^{C_g}$, then $B\ll X^{1/3}L^{C_g/3}$.
Since the product of all three scales is comparable to $X$ and each
is at most $B$, each is also at least $X^{1/3}$ times a fixed negative
power of $L$.  Thus only $O((1+\log L)^3)$ triples of dyads occur.
Their roles in \eqref{eq:distinguished-tuples} and their orderings
have bounded multiplicity.  The prime ideal theorem, or its upper
bound alone, gives \eqref{eq:decomposition-triple-moments}.
This proves the exceptional assertions of the proposition in this branch.

\subsection{Stopping without coupled coefficients}

It remains to treat $R<X^{s_0}$.  Exact Mobius expansion gives
\begin{equation}\label{eq:rough-mobius-expansion}
 \mathcal F_w(h)=
 \sum_{mu=h}\mu(m)\prod_{\pi\mid m}\psi_w(\pi).
\end{equation}
The $m$ primes have norm at most $2w$, whereas $u$ is unrestricted.
Partition their possible prime norms into $O(L)$ bins, in descending
order, each with upper/lower ratio at most $1+\delta_0$.  Here
$\delta_0>0$ is fixed sufficiently small in terms of $\kappa$.
On ordinary blocks include $L^{C_r}$ as a boundary.  Endpoints can
be assigned consistently to either adjacent bin; use the lower side
at this particular boundary.

If a prime lies in a bin of lower endpoint $\ell$, replace its norm
by $\ell$ when forming a \emph{surrogate product}, denoted $S$.
There are $O(L)$ factors in any nonzero term on the envelope.
Consequently $\delta_0$ can be chosen so that
\begin{equation}\label{eq:surrogate-norm}
 S(m')\le N(m')\le X^{\kappa/4}S(m')
\end{equation}
for every selected subproduct $m'$.
Start the running surrogate at the exact value $N(r)$.

On ordinary blocks process first the bins above $L^{C_r}$ and stop
on first reaching $X^{.36}$.  If these bins are exhausted without
a stop, process the remaining bins, now stopping on first reaching
$X^{.38}$.  On upper blocks process all bins with the single stopping
threshold $X^{1/3-\kappa}$.  Since $N(r)\asymp R<X^{s_0}$, the
initial value lies below every applicable first threshold for large $X$.

If there is no stop, then $N(r)S(m)<Z$, where the final threshold
$Z$ is $X^{.38}$ on ordinary blocks and $X^{1/3-\kappa}$ on upper
blocks.  All predicates involve $r,m$ alone.  Grouping
by dyads of $N(rm)$ therefore gives Type~I coefficients, with $u$
unrestricted.  By \eqref{eq:surrogate-norm},
$N(rm)\le X^{\kappa/4}N(r)S(m)<X^{\kappa/4}Z$.
Thus the outer norm is less
than $X^{.38+\kappa/4}<X^{.40}$ on ordinary blocks, and less than
$X^{1/3-3\kappa/4}<X^{1/3-\kappa/2}$ on upper blocks.
For outer variable $v=rm$ on a dyad of length $R'$ the inner weight is
$V(N(v)U x/X)$, where $U=X/R'$.  Its derivatives are bounded uniformly
in $v$, as required in Propositions~\ref{prop:typeI} and
\ref{prop:typeIheight}.

For stopped terms write $m=m_+m_-$ and set
\[
 b=rm_+,\qquad a=m_-u.
\]
The crossing prime has norm at most $2X^\xi$.  Thus ordinary stopped
terms have
\[
 X^{.36}\le N(b)\ll X^{.38+\xi+\kappa/4}<X^{.40},
\]
and upper stopped terms have
\[
 X^{1/3-\kappa}\le N(b)
 \ll X^{1/3-\kappa+\xi+\kappa/4}<X^{.40}.
\]
Sharp norm dyads may now be inserted separately on $a,b$; prime-polynomial
structure is not required in this range.

\begin{lemma}[Exact separation at the stopping bin]
\label{lem:stopping-separation}
The stopped terms are exact sums with independent coefficients on
$a$ and $b$.  The number of additional outer indices is a fixed
power of $L$, and the coefficients have fixed divisor bounds
independent of $C_r$.
\end{lemma}

\begin{proof}
Fix the last bin, its lower endpoint $\ell$, and the numbers $k_+\ge1$
and $k_-\ge0$ of selected and remaining factors in this bin.
All factors in earlier bins belong to $m_+$ and all factors in later
bins to $m_-$.  If $Z$ is the applicable threshold, the crossing
predicate is
\begin{equation}\label{eq:crossing-predicate}
 N(r)S(m_+)/\ell<Z\le N(r)S(m_+).
\end{equation}
For an ordinary stop below the roughness boundary, additionally require
$N(r)S(m_+^{>})<X^{.36}$, where $m_+^{>}$ contains the factors above
that boundary.  Every factor in those earlier bins is already on
the selected side.  Hence both predicates depend only on the
$b$-side variables.

The surrogate uses the same $\ell$ for every prime in the final bin.
For a fixed original squarefree $m$, exactly
$\binom{k_++k_-}{k_+}$ selections have the stated counts, and either
all or none satisfy the crossing test.  Give each selection the scalar
\begin{equation}\label{eq:stopping-scalar}
 \binom{k_++k_-}{k_+}^{-1}.
\end{equation}
Their sum restores the original coefficient exactly.  The Mobius
factor and all prime weights split multiplicatively between $m_+$
and $m_-$.  Once the last bin and its two counts have been fixed,
the $b$ coefficient sums only over $r,m_+$, and the $a$ coefficient
only over $m_-,u$.

More explicitly, let $v(r)$ denote the weighted distinguished tuple
sum, including its factor $1/i!$, and put
$q(d)=\mu(d)\prod_{\pi\mid d}\psi_w(\pi)$.
If $j$ is the last bin and $\mathcal C_{j,k}(r,d)$ is the crossing
condition with exactly $k$ factors of $d$ in bin $j$ and no factors
in later bins, define
\[
 \beta_{j,k}(b)=\sum_{rd=b}v(r)q(d)
                    \boldsymbol1_{\mathcal C_{j,k}(r,d)},
 \qquad
 \alpha_{j,l}(a)=\sum_{eu=a}q(e)
   \boldsymbol1_{\substack{e\text{ has factors only in bins }j,j+1,\ldots\\
                           e\text{ has exactly }l\text{ factors in bin }j}}.
\]
Then the stopped sum against any squarefree-supported kernel $K$ is
exactly
\[
 \sum_{j,k\ge1,l\ge0}\binom{k+l}{k}^{-1}
       \sum_{a,b}\alpha_{j,l}(a)\beta_{j,k}(b)K(ab).
\]
Here $K$ includes the product envelope, and the notation
$\mathcal C_{j,k}$ includes the failed first-stage test when needed.

One may enlarge these independent sums by allowing overlaps among
$r,m_+,m_-,u$.  Each added tuple has nonsquarefree total product and
is individually annihilated by both kernels, before coefficient
collection.  After collection, $a,b$ can be restricted separately
squarefree; a remaining cross-side common prime is killed by
$G(ab)$ and by $\mu^2(ab)$.  Thus independence does not require an
unaccounted coprimality restriction.

The bin groups are unordered divisors.  Only the already bounded
distinguished tuple is ordered.  Their multiplicities are therefore
bounded by fixed powers of $\tau(a),\tau(b)$, and the scalar
\eqref{eq:stopping-scalar} is at most $1$.  There are $O(L)$ choices
for the final bin and $O(L)$ for each of its two counts.  The complete
bin occupancy vector is not an outer index: all other bin conditions
remain inside the respective coefficient.  Adding the boundary
$L^{C_r}$ adds at most one bin and changes none of these powers.
\end{proof}

\subsection{Sparse late stops}

For an ordinary stop using a prime of norm at most $L^{C_r}$,
the high-bin running surrogate is less than $X^{.36}$, while the
final surrogate is at least $X^{.38}$.  Thus $b$ has a divisor $d$
with
\[
 N(d)\ge X^{.02},\qquad
 N(\pi)\le L^{C_r}\quad(\pi\mid d).
\]
For fixed $C_r$, choose a fixed $\sigma>0$ with $C_r\sigma<1$.
Rankin's inequality gives
\begin{align}
 \sum_{\substack{N(d)\ge X^{.02}\\
                  N(\pi)\le L^{C_r}\ (\pi\mid d)}}\frac1{N(d)}
 &\le X^{-.02\sigma}
 \prod_{N(\pi)\le L^{C_r}}
       (1-N(\pi)^{-1+\sigma})^{-1}
 \label{eq:late-stop-rankin}\\
 &\le X^{-.02\sigma+o(1)}.\notag
\end{align}
Indeed the logarithm of the product is $O(L^{C_r\sigma})=o(L)$,
using the prime ideal counting upper bound.  The number of multiples
of $d$ on a dyad of length $B$ is $O(B/N(d))$ when $N(d)\ll B$;
it follows that the allowed $b$ have cardinality $O(BX^{-\delta})$
for some $\delta>0$.  Fixed divisor bounds preserve a smaller positive
power saving in $\sum_b|\beta_b|^2$.

Here $X^{.35}\ll B<X^{.40}$ and $A\asymp X/B$, so
$A+B+(AB)^{2/3}\ll X^{2/3}$.  Cauchy's inequality and the cubic
large sieve yield
\[
 \left|\sum_{a,b}\alpha_a\beta_bG(ab)N(ab)^{it}
                 V(N(ab)/X)\right|
 \ll X^{5/6-\eta}
\]
for a fixed $\eta>0$.  The envelope separates by smooth Mellin
inversion.  The model has the same saving by direct sparse counting,
since $N(ab)^{-1/6}\asymp X^{-1/6}$.  Norm twists have modulus $1$.
This proves part~(5) of Proposition~\ref{prop:decomposition}, even
after summing the polynomially many pieces.

\subsection{Coefficient moments and the Siegel--Walfisz condition}

All coefficients constructed above satisfy fixed divisor bounds.
For completeness, these give the requisite moment estimates also
after collecting equal norms.  The number of ideals of norm $m$ is
at most the ordinary divisor function $d(m)$, and the ideal divisor
function of any such ideal is at most $d(m)^2$.  Each fixed power
of $d(m)$ is bounded by a higher divisor function, whose sum up to
$Y$ is $O(Y(\log(2Y))^C)$.  Thus every fixed coefficient moment
on a dyad has the claimed bound.  This reasoning also bounds the
total absolute multiplicity of complementary factorizations below.

The nonsparse stopped ordinary coefficients are rough: distinguished
primes have norm at least $w$, and selected bin primes exceed
$L^{C_r}$.  Pure prime-tuple coefficients are rough as well, since
all factors exceed $w$.

\begin{lemma}[The sequence condition]\label{lem:decomposition-sw}
Every nonsparse ordinary $b$ coefficient satisfies \textup{(SW)}:
against any nonprincipal cubic Hecke character of the indicated
trivial infinite parameter and primitive conductor of fixed
polylogarithmic norm, its sum is $O_D(BL^{-D})$ for every fixed $D$.
This remains true with the permitted polynomial-norm coprimality
exclusions and fixed-polylogarithmic norm twists.
\end{lemma}

\begin{proof}
First discard terms all of whose prime factors have norm at most
$y=\exp(\sqrt L)$.  Apply Rankin's inequality with
$\sigma=1/\sqrt L$.  Fixed divisor multiplicities can be included in
the Euler product; its logarithm is $O(\log L)$, since
\[
 \sum_{N(\pi)\le y}N(\pi)^{-1+1/\sqrt L}=O(\log L).
\]
As $B$ is a fixed positive power of $X$, the resulting total absolute
weight is at most $B\exp(-c\sqrt L)L^C$.  This gives every desired
logarithmic saving.

In the remaining terms select a largest prime, breaking ties by a
fixed order.  Fix its role, its norm dyad $P$, and all other factors.
There are boundedly many distinguished roles and one unordered
selected-bin role; no permutation of a bin product is introduced.
We have $P\gg\exp(\sqrt L)$.

Refine its range by the bins.  If the variable prime lies in $m_+$,
all surrogate data are constant within its bin.  If it lies in $r$,
then $N(r)$ is a fixed multiple of its norm.  Accordingly,
\eqref{eq:crossing-predicate}, the possible failed first-stage test,
the final $b$ dyad, and the largest-prime condition each give an
interval restriction in its norm.  Their intersection has only
polynomially many interval pieces as the bin varies.  After the
other factors are fixed, the Mobius sign is constant, and the remaining
prime weights are smooth with bounded scaled derivatives.

The prime ideal Siegel--Walfisz estimate from the background inputs
applies on this prime range.  A conductor bound $L^C$ is a bound
$(\log P)^{2C+O(1)}$, and arbitrary powers of $\log P$ give arbitrary
powers of $L$ in the saving.  Partial summation absorbs the smooth
weights and any fixed-polylogarithmic norm twist.  It applies to
arbitrary subintervals by subtracting two initial-interval estimates.

Tie decisions and distinctness cause no loss: in this quadratic field
there are at most two prime ideals of any prime norm, and at most one
of an inert prime-square norm.  Once the other factors are fixed,
tie rules affect only endpoint norms.  Coprimality to a number of
polynomial norm removes $O(L)$ additional primes.  These deletions
are negligible relative to $P L^{-D}$, since
$P\gg\exp(\sqrt L)$.

For a fixed dyad $P$, the total absolute weight of complementary
factorizations is at most $(B/P)L^{O(1)}$, by the fixed divisor
multiplicity bounds.  Summing the prime cancellation over these
choices, over roles, and over the polynomially many dyads and interval
pieces proves the assertion.  The argument for pure prime-tuple
coefficients is the same, without the surrogate restrictions.
\end{proof}

\begin{proof}[Completion of Proposition~\ref{prop:decomposition}]
The identities \eqref{eq:prime-detector},
\eqref{eq:distinguished-tuples}, and \eqref{eq:rough-mobius-expansion}
are exact on squarefree numbers.  Lemma~\ref{lem:prime-grouping}
treats the large-$R$ branch, and Lemma~\ref{lem:stopping-separation}
treats the small-$R$ branch.  The actual norm ranges proved above
give parts~(1) and~(2).  Lemma~\ref{lem:decomposition-sw} gives
part~(3), and the three-prime counting argument gives part~(4).
Ordinary stopped terms have $H\gg X^{.08}$, so no additional
small-$H$ configuration is omitted.  The sparse estimate gives
part~(5).

The total number of pieces is a fixed power of $L$: the distinguished
tuple has bounded length, the bin and endpoint-count indices have
polynomially many choices, and the side dyads add only polynomial
factors.  No full occupancy vector of the bins is enumerated.
The exponents in this complexity and in the divisor bounds are
independent of $C_r,C_g$; those parameters affect only fixed constants
in the exceptional triple count, not its polynomial dependence on
$\log L$.

The parameter choices are therefore noncircular.  First fix the
absolute power-width constant for $B\le X^{.40}$ in
Proposition~\ref{prop:dispersion}; then fix $\kappa,\rho,\xi$ and
the resulting bounded tuple complexity.  Choose the logarithmic
output and localization savings, and the large-height threshold
power, after that complexity.  Finally choose $C_g,C_r$ and the
needed SW saving powers.  The large-$B$ prime-polynomial estimates
may depend on the now fixed tuple complexity and $\xi$; they do not
determine the earlier absolute power-width constant.
\end{proof}
